% Part: computability
% Chapter: computability-theory
% Section: complement-ce

\documentclass[../../../include/open-logic-section]{subfiles}

\begin{document}

\olfileid[pt]{cmp}{thy}{cmp}
\olsection{Conjuntos computavelmente enumeráveis não são fechados sob
complementação}

Suponha que $A$ seja computavelmente enumerável. O complemento de~$A$,
$\Complement{A} = \Nat \setminus A$, também é sempre computavelmente
enumerável? O teorema e o corolário a seguir mostram que a resposta é
``não''.

\begin{thm}
\ollabel{thm:ce-comp}
Seja $A$ um conjunto qualquer de números naturais. Então $A$ é
computável se, e somente se, tanto $A$ quanto $\Complement{A}$ são
computavelmente enumeráveis.
\end{thm}

\begin{proof}
O sentido direto é fácil: se $A$ é computável, então
$\Complement{A}$ também é computável ($\Char{A} = 1 \tsub
\Char{\Complement{A}}$) e, portanto, ambos são computavelmente
enumeráveis.

No outro sentido, suponha que $A$ e~$\Complement{A}$ sejam ambos
computavelmente enumeráveis. Seja $A$ o domínio de~$\cfind{d}$, e seja
$\Complement{A}$ o domínio de~$\cfind{e}$. Defina $h$ por
\[
h(x) = \umin{s}{(T(d,x,s) \lor T(e,x,s))}.
\]
Em outras palavras, na entrada~$x$, $h$~procura uma computação que pare
de~$\cfind{d}$ ou uma computação que pare de~$\cfind{e}$. Se
$x \in A$, ela terá êxito no primeiro caso; se
$x \in \Complement{A}$, terá êxito no segundo. Logo, $h$~é uma função
computável total. Mas agora temos que, para todo~$x$, $x \in A$ se, e
somente se, $T(e, x, h(x))$, isto é, se $\cfind{e}$ é a função que
está definida. Como $T(e, x, h(x))$ é uma relação computável,
$A$~é computável.
\end{proof}

\begin{explain}
É mais fácil entender o que ocorre em termos computacionais informais:
para decidir $A$, na entrada $x$, procure computações que parem de
$\cfind{e}$ e $\cfind{f}$. Uma delas necessariamente parará; se for
$\cfind{e}$, então $x$ pertence a~$A$; caso contrário, $x$ pertence
a~$\Complement{A}$.
\end{explain}

\begin{cor}
\ollabel{cor:comp-k}
$\Complement{K_0}$ não é computavelmente enumerável.
\end{cor}

\begin{proof}
Sabemos que $K_0$ é computavelmente enumerável, mas não é computável.
Se $\Complement{K_0}$ fosse computavelmente enumerável, então $K_0$
seria computável por \olref{thm:ce-comp}, contradizendo
\olref[ncp]{thm:K-0}.
\end{proof}

\end{document}
